% Options for packages loaded elsewhere
\PassOptionsToPackage{unicode}{hyperref}
\PassOptionsToPackage{hyphens}{url}
\documentclass[
  ignorenonframetext,
]{beamer}
\newif\ifbibliography
\usepackage{pgfpages}
\setbeamertemplate{caption}[numbered]
\setbeamertemplate{caption label separator}{: }
\setbeamercolor{caption name}{fg=normal text.fg}
\beamertemplatenavigationsymbolsempty
% remove section numbering
\setbeamertemplate{part page}{
  \centering
  \begin{beamercolorbox}[sep=16pt,center]{part title}
    \usebeamerfont{part title}\insertpart\par
  \end{beamercolorbox}
}
\setbeamertemplate{section page}{
  \centering
  \begin{beamercolorbox}[sep=12pt,center]{section title}
    \usebeamerfont{section title}\insertsection\par
  \end{beamercolorbox}
}
\setbeamertemplate{subsection page}{
  \centering
  \begin{beamercolorbox}[sep=8pt,center]{subsection title}
    \usebeamerfont{subsection title}\insertsubsection\par
  \end{beamercolorbox}
}
% Prevent slide breaks in the middle of a paragraph
\widowpenalties 1 10000
\raggedbottom
\AtBeginPart{
  \frame{\partpage}
}
\AtBeginSection{
  \ifbibliography
  \else
    \frame{\sectionpage}
  \fi
}
\AtBeginSubsection{
  \frame{\subsectionpage}
}
\usepackage{iftex}
\ifPDFTeX
  \usepackage[T1]{fontenc}
  \usepackage[utf8]{inputenc}
  \usepackage{textcomp} % provide euro and other symbols
\else % if luatex or xetex
  \usepackage{unicode-math} % this also loads fontspec
  \defaultfontfeatures{Scale=MatchLowercase}
  \defaultfontfeatures[\rmfamily]{Ligatures=TeX,Scale=1}
\fi
\usepackage{lmodern}
\ifPDFTeX\else
  % xetex/luatex font selection
\fi
% Use upquote if available, for straight quotes in verbatim environments
\IfFileExists{upquote.sty}{\usepackage{upquote}}{}
\IfFileExists{microtype.sty}{% use microtype if available
  \usepackage[]{microtype}
  \UseMicrotypeSet[protrusion]{basicmath} % disable protrusion for tt fonts
}{}
\makeatletter
\@ifundefined{KOMAClassName}{% if non-KOMA class
  \IfFileExists{parskip.sty}{%
    \usepackage{parskip}
  }{% else
    \setlength{\parindent}{0pt}
    \setlength{\parskip}{6pt plus 2pt minus 1pt}}
}{% if KOMA class
  \KOMAoptions{parskip=half}}
\makeatother
\usepackage{longtable,booktabs,array}
\newcounter{none} % for unnumbered tables
\usepackage{calc} % for calculating minipage widths
\usepackage{caption}
% Make caption package work with longtable
\makeatletter
\def\fnum@table{\tablename~\thetable}
\makeatother
\setlength{\emergencystretch}{3em} % prevent overfull lines
\providecommand{\tightlist}{%
  \setlength{\itemsep}{0pt}\setlength{\parskip}{0pt}}
\usepackage{bookmark}
\IfFileExists{xurl.sty}{\usepackage{xurl}}{} % add URL line breaks if available
\urlstyle{same}
\hypersetup{
  hidelinks,
  pdfcreator={LaTeX via pandoc}}

\author{\texorpdfstring{}{}}
\date{}

\begin{document}

\begin{frame}[fragile]{Methodology: Pipeline Design and Formal
Definitions \label{sec:methods}}
\protect\phantomsection\label{methodology-pipeline-design-and-formal-definitions}
This section describes the five-stage computational meta-analysis
pipeline. Each stage corresponds to a tested, independently executable
script that reads upstream outputs and produces structured artifacts.
The pipeline extends the systematic literature analysis approach of
Knight et al.~\citep{knight2022fep}---which combined manual annotation
with ontology-based automated analysis---by substituting manual coding
with fully automated, LLM-driven assertion extraction and
citation-weighted hypothesis scoring.

\begin{block}{Pipeline Overview}
\protect\phantomsection\label{pipeline-overview}
{\def\LTcaptype{none} % do not increment counter
\begin{longtable}[]{@{}
  >{\raggedright\arraybackslash}p{(\linewidth - 8\tabcolsep) * \real{0.2000}}
  >{\raggedright\arraybackslash}p{(\linewidth - 8\tabcolsep) * \real{0.2000}}
  >{\raggedright\arraybackslash}p{(\linewidth - 8\tabcolsep) * \real{0.2000}}
  >{\raggedright\arraybackslash}p{(\linewidth - 8\tabcolsep) * \real{0.2000}}
  >{\raggedright\arraybackslash}p{(\linewidth - 8\tabcolsep) * \real{0.2000}}@{}}
\toprule\noalign{}
\begin{minipage}[b]{\linewidth}\raggedright
Stage
\end{minipage} & \begin{minipage}[b]{\linewidth}\raggedright
Script
\end{minipage} & \begin{minipage}[b]{\linewidth}\raggedright
Primary Input
\end{minipage} & \begin{minipage}[b]{\linewidth}\raggedright
Primary Output
\end{minipage} & \begin{minipage}[b]{\linewidth}\raggedright
Section
\end{minipage} \\
\midrule\noalign{}
\endhead
1 & \texttt{01\_literature\_search.py} & API queries &
\texttt{corpus.jsonl} & \hyperref[sec:methods_retrieval]{Retrieval} \\
2 & \texttt{02\_meta\_analysis\_pipeline.py} & \texttt{corpus.jsonl} &
Classification, temporal, TF-IDF, NMF, citation network JSONs &
\hyperref[sec:methods_bibliometrics]{Bibliometrics} \\
3 & \texttt{03\_build\_knowledge\_graph.py} & \texttt{corpus.jsonl} &
\texttt{nanopublications.jsonl}, \texttt{nanopublications.trig}, scores
& \hyperref[sec:methods_kg]{Knowledge Graph} \\
4 & \texttt{04\_generate\_figures.py} & All Stage 2--3 JSONs & 16
publication-ready PNGs & \hyperref[sec:methods_viz]{Visualization} \\
5 & \texttt{05\_inject\_variables.py} & All output JSONs & Rendered
manuscript Markdown & \hyperref[sec:methods_viz]{Injection} \\
\bottomrule\noalign{}
\end{longtable}
}

Scripts act as thin orchestrators that import methods from tested
library modules and handle file I/O. All computation resides in the
\texttt{src/} packages; no analysis logic is embedded in scripts.
End-to-end pipeline execution completes in under one hour on commodity
hardware (excluding LLM extraction, which depends on model size and
inference backend); all stochastic components use fixed random seeds for
deterministic reproduction.
\end{block}
\end{frame}

\end{document}
